% Ad Atticum 13.32 --- headnote
% ad-atticum-13-32, 29 May 45 BC, Tusculanum

\textit{Cicero to Atticus, written at the Tusculan villa on 29 May
45 BC --- Perseus dateline \textit{Scr. in Tusculano iv K. Iun. a.
709 (45)}. The last of the late-May daily letters from Tusculum.
Three short sections: the Faberian negotiations are still the
hinge on which the gardens purchase --- and behind the gardens
purchase, the shrine for Tullia --- depends; Cicero needs more
Dicaearchus for the Academica revision; and the Mummius-commission
question raised in 13.30 is now resolved (Postumius, not
Tuditanus, supplied by Atticus's recollection of a statue at the
Isthmus).}

\textit{Five Greek phrases punctuate the page. \textit{peri psychēs}
``On the Soul,'' \textit{katabaseōs} ``of the \textit{Descent},''
and \textit{Tripolitikon} are book titles --- Dicaearchus's
philosophical and political dialogues that Cicero wants on his
desk for the new four-book Academica. \textit{dia sēmeiōn},
``in shorthand,'' apologises for an obscure earlier note (Cicero
had been writing in coded brachygraphy and so Atticus had missed
the point); \textit{syllogon} and \textit{pompeusai kai tois
prosōpois} pick up the metaphor running through the cluster ---
the dialogue is being cast like a Greek procession, with the right
\textit{prosōpa}, ``faces'' or ``characters,'' marching in their
proper places. The reference to ``new prefaces'' in which Catulus
and Lucullus are each praised is to the still-extant earlier
two-book version of the Academica (the \textit{Catulus} and
\textit{Lucullus}), which Cicero is now in the process of
superseding by the four-book Varro-dedication. The whole cluster
of 13.29--32 belongs to that pivot.}
